\documentclass[12pt]{article}
\usepackage[T1]{fontenc}
\usepackage[utf8]{inputenc}
\usepackage{lmodern}
\usepackage{textcomp}
\usepackage{enumitem}
\usepackage{geometry}
\geometry{margin=1in}
\begin{document}
\begin{center}
Answer Key - Physics - Undergraduate Level Assessment 19 \
\small (Answers are ordered exactly as in the question paper.)
\end{center}

\section*{Multiple Choice}
\begin{enumerate}[label=\arabic*.]
\item \textbf{Answer:} B
\\ \textit{Options:} A) an average of 10 times, with an rms deviation of about 4; B) an average of 10 times, with an rms deviation of about 3; C) an average of 10 times, with an rms deviation of about 1; D) an average of 10 times, with an rms deviation of about 0.1
\\ \textit{Note:} The quantum efficiency of 0.1 means that on average 10\% of the incoming photons are detected, so with 100 photons sent in, the expected number detected is 10, while the rms deviation, which quantifies the statistical variation in this Poisson process, is approximately the square root of the average detected (√10), giving about 3.
\item \textbf{Answer:} B
\\ \textit{Options:} A) 2; B) 250; C) 5,000; D) 10,000
\\ \textit{Note:} The resolving power of a spectrometer is defined as the ratio of the wavelength (λ) to the minimum difference in wavelength (Δλ) that can be resolved, so for wavelengths of 500 nm and 502 nm, the resolving power is calculated as 500 nm / 2 nm = 250.
\item \textbf{Answer:} A
\\ \textit{Options:} A) decreased by a factor of 9; B) decreased by a factor of 49; C) decreased by a factor of 81; D) increased by a factor of 9
\\ \textit{Note:} The emission spectrum of Li++ is identical to that of hydrogen because both systems are hydrogen-like, and since Li++ has a nuclear charge (Z) of 3, the wavelengths of its spectral lines are scaled by the factor of 1/Z², resulting in wavelengths decreased by a factor of 9 (1/3²) compared to hydrogen.
\item \textbf{Answer:} D
\\ \textit{Options:} A) 0.1 GeV/$c^2$; B) 0.2 GeV/$c^2$; C) 0.5 GeV/$c^2$; D) 1.0 GeV/$c^2$
\\ \textit{Note:} The correct answer is derived from the relativistic energy-momentum relation \(E^2 = (pc)^2 + (m_0 c^2)^2\), where substituting the given values for total energy and momentum allows us to solve for the rest mass, yielding approximately 1.0 GeV/c².
\item \textbf{Answer:} D
\\ \textit{Options:} A) Specific heat; B) Thermal conductivity; C) Electrical resistivity; D) Hall coefficient
\\ \textit{Note:} The sign of the charge carriers in a doped semiconductor can be determined from the Hall coefficient, as it indicates whether the dominant charge carriers are positive (holes) or negative (electrons) based on the direction of the induced voltage in the presence of a magnetic field.
\end{enumerate}

\section*{True / False}
\begin{enumerate}[label=\arabic*.]
\setcounter{enumi}{5}
\item \textbf{Answer:} True
\\ \textit{Options:} A) True; B) False
\\ \textit{Note:} In a constant volume process, the work done on or by the system is zero, so the change in internal energy of an ideal gas is equal to the heat added, according to the first law of thermodynamics.
\item \textbf{Answer:} True
\\ \textit{Options:} A) True; B) False
\\ \textit{Note:} In the diamond structure of elemental carbon, each carbon atom forms four covalent bonds with four other carbon atoms located at the corners of a tetrahedron, making them the nearest neighbors.
\item \textbf{Answer:} True
\\ \textit{Options:} A) True; B) False
\\ \textit{Note:} The total spin quantum number of the electrons in the ground state of neutral nitrogen is 3/2 because nitrogen has three unpaired electrons, each contributing a spin of 1/2, leading to a total spin of 3/2.
\item \textbf{Answer:} False
\\ \textit{Options:} A) True; B) False
\\ \textit{Note:} The statement is false because bosons have symmetric wave functions and do not obey the Pauli exclusion principle, which applies only to fermions.
\item \textbf{Answer:} True
\\ \textit{Options:} A) True; B) False
\\ \textit{Note:} An atom with filled n = 1 and n = 2 energy levels contains a total of 2 electrons in n = 1 and 8 electrons in n = 2, totaling 10 electrons, which corresponds to the electron configuration of the noble gas neon.
\end{enumerate}

\section*{Long Form Response}
\begin{enumerate}[label=\arabic*.]
\setcounter{enumi}{10}
\item \textbf{Answer:} To calculate the distance a particle travels before decaying in the lab frame, we first need to consider the effects of time dilation from special relativity. The proper decay time of the particle is 2.0 ms, but due to its velocity of 0.60 c, the time observed in the lab frame will be longer. The time dilation factor, or Lorentz factor (γ), is calculated as γ = 1 / \ensuremath{\sqrt}(1 - ($v^2$/$c^2$)), which gives γ \ensuremath{\approx} 1.25. Therefore, the dilated decay time in the lab frame is 2.0 ms \ensuremath{\times} 1.25 = 2.5 ms. The distance covered before decaying is then calculated using the formula distance = speed \ensuremath{\times} time. Substituting in the values, we find that the distance is 0.60 c \ensuremath{\times} 2.5 ms = 450 m. Thus, the particle covers a distance of 450 meters before decaying, illustrating the impact of relativistic effects on its observed lifetime.
\\ \textit{Note:} correct because it applies the concept of time dilation from special relativity to determine the observed decay time in the lab frame, allowing for the calculation of the distance traveled by the particle before decaying at a velocity of 0.60 c.
\item \textbf{Answer:} The relationship between the absolute temperature of a blackbody and the energy it radiates per second per unit area is described by Stefan-Boltzmann Law, which states that the total energy radiated is proportional to the fourth power of the blackbody's absolute temperature (T). Specifically, this means that when the temperature of a blackbody increases, the energy output increases significantly; for example, if the temperature is increased by a factor of 3, the energy radiated per second per unit area increases by a factor of $3^4$, or 81. This dramatic increase occurs because higher temperatures enhance the kinetic energy of particles, which in turn leads to more energetic electromagnetic radiation being emitted. Thus, the relationship illustrates how sensitive blackbody radiation is to temperature changes, highlighting the principle that even small increases in temperature yield substantial increases in energy output.
\\ \textit{Note:} correct because it accurately describes the Stefan-Boltzmann Law, which quantitatively relates the absolute temperature of a blackbody to its energy output, illustrating that energy radiated increases dramatically with temperature, specifically by the fourth power of the temperature increase.
\end{enumerate}

\vfill
\noindent\textit{End of Answer Key}
\end{document}